% GENERATED FILE. Edit canonical lessons, then run npm run generate:books.
% canonical-source-hash: fnv1a64:37e3fcf7b420e3c9
% canonical-lessons: TA-C499-varatatcanai, TA-C499-acirvatam, TA-C499-mana, TA-C499-tampati, TA-C499-illaram

\chapter{Getting Married}
\label{ch:ta-getting-married}
\hlchaptermodality{499}

\begin{chapteropening}
\textbf{By the end of this chapter:} \emph{I can talk about dowry, blessings, marrying and married couples.}

Complete the last lesson of chapter 499: I can talk about dowry, blessings, marrying and married couples.
\end{chapteropening}

\section[\texorpdfstring{\ta{வரதட்சணை} (varataṭcaṇai)}{வரதட்சணை (varataṭcaṇai)}]{\ta{வரதட்சணை} (varataṭcaṇai) — a dowry}

\label{lesson:TA-C499-varatatcanai}



\begin{quote}
Before the new one: say the Tamil for a prospective groom, then the Tamil for consent.
\end{quote}

\subsection*{You'll want to know: \ta{வரதட்சணை}}
\textbf{\ta{வரதட்சணை}} — \emph{varataṭcaṇai} — \textquotedblleft{}a dowry\textquotedblright{}.

\begin{quote}
\ta{அவர்கள்} \ta{வரதட்சணை} \ta{வாங்காமல்} \ta{திருமணம்} \ta{செய்தார்கள்}.
\end{quote}

\begin{quote}
\emph{avarkaḷ varataṭcaṇai vāṅkāmal tirumaṇam ceytārkaḷ.}
\end{quote}

\begin{quote}
They married without taking a dowry.
\end{quote}

\ta{வரதட்சணை} \ta{வாங்கு}, to take a dowry.

\begin{cousinweb}
\textbf{\ta{வரதட்சணை}} (\emph{varataṭcaṇai}) is borrowed from Sanskrit, built from \emph{vara}, \textquotedblleft{}bridegroom\textquotedblright{} (CDIAL 11309), and \emph{dakṣiṇā}, \textquotedblleft{}a gift\textquotedblright{} (CDIAL 6120); Tamil writes \emph{kṣ} as \ta{ட்ச}.
\end{cousinweb}

\begin{grammarlens}[title={What you've built}]
One more word for describing what you lived and felt.
\end{grammarlens}

\subsection*{Your turn}
\begin{itemize}
\raggedright
  \item \emph{Say it:} \emph{varataṭcaṇai}
  \item \emph{Say it:} \emph{varataṭcaṇai}, once more
  \item \emph{Say it:} \emph{avarkaḷ varataṭcaṇai vāṅkāmal tirumaṇam ceytārkaḷ.} — They married without taking a dowry.
  \item \emph{Recall:} say the Tamil for a prospective groom, then the Tamil for consent, then say \emph{varataṭcaṇai} again
\end{itemize}

\begin{tcolorbox}[breakable,colback=teal!4,colframe=teal!35!black,title={Before you move on}]
What does \ta{வரதட்சணை} mean? (\textquotedblleft{}A dowry\textquotedblright{}.) Say it once more.
\end{tcolorbox}
\section[\texorpdfstring{\ta{ஆசீர்வாதம்} (ācīrvātam)}{ஆசீர்வாதம் (ācīrvātam)}]{\ta{ஆசீர்வாதம்} (ācīrvātam) — a blessing}

\label{lesson:TA-C499-acirvatam}



\begin{quote}
Before the new one: say the Tamil for consent, then the Tamil for a dowry.
\end{quote}

\subsection*{You'll want to know: \ta{ஆசீர்வாதம்}}
\textbf{\ta{ஆசீர்வாதம்}} — \emph{ācīrvātam} — \textquotedblleft{}a blessing\textquotedblright{}.

\begin{quote}
\ta{திருமணத்துக்கு} \ta{முன்} \ta{அவர்கள்} \ta{பெரியவர்களிடம்} \ta{ஆசீர்வாதம்} \ta{வாங்கினார்கள்}.
\end{quote}

\begin{quote}
\emph{tirumaṇattukku muṉ avarkaḷ periyavarkaḷiṭam ācīrvātam vāṅkiṉārkaḷ.}
\end{quote}

\begin{quote}
Before the wedding they received the elders’ blessing.
\end{quote}

\ta{ஆசீர்வாதம்} \ta{வாங்கு}, to receive a blessing.

\begin{cousinweb}
\textbf{\ta{ஆசீர்வாதம்}} (\emph{ācīrvātam}) is borrowed from Sanskrit \emph{āśīrvāda}, \textquotedblleft{}benediction\textquotedblright{} (Monier-Williams).
\end{cousinweb}

\begin{grammarlens}[title={What you've built}]
One more word for describing what you lived and felt.
\end{grammarlens}

\subsection*{Your turn}
\begin{itemize}
\raggedright
  \item \emph{Say it:} \emph{ācīrvātam}
  \item \emph{Say it:} \emph{ācīrvātam}, once more
  \item \emph{Say it:} \emph{tirumaṇattukku muṉ avarkaḷ periyavarkaḷiṭam ācīrvātam vāṅkiṉārkaḷ.} — Before the wedding they received the elders’ blessing.
  \item \emph{Recall:} say the Tamil for consent, then the Tamil for a dowry, then say \emph{ācīrvātam} again
\end{itemize}

\begin{tcolorbox}[breakable,colback=teal!4,colframe=teal!35!black,title={Before you move on}]
What does \ta{ஆசீர்வாதம்} mean? (\textquotedblleft{}A blessing\textquotedblright{}.) Say it once more.
\end{tcolorbox}
\section[\texorpdfstring{\ta{மண} (maṇa)}{மண (maṇa)}]{\ta{மண} (maṇa) — to marry}

\label{lesson:TA-C499-mana}



\begin{quote}
Before the new one: say the Tamil for a dowry, then the Tamil for a blessing.
\end{quote}

\subsection*{You'll want to know: \ta{மண}}
\textbf{\ta{மண}} — \emph{maṇa} — \textquotedblleft{}to marry\textquotedblright{}.

\begin{quote}
\ta{அவர்} \ta{ஒரு} \ta{மருத்துவரை} \ta{மணந்தார்}.
\end{quote}

\begin{quote}
\emph{avar oru maruttuvarai maṇantār.}
\end{quote}

\begin{quote}
He married a doctor.
\end{quote}

Past (I …): \ta{மணந்தேன்}, \emph{maṇantēṉ}.

Formal and literary; the everyday phrase is \ta{திருமணம்} \ta{செய்}. The one married takes -\ta{ஐ}.

\begin{cousinweb}
\textbf{\ta{மண}} (\emph{maṇa}) is native Dravidian (DEDR 4667). DEDR keeps \ta{மண} \textquotedblleft{}to emit fragrance\textquotedblright{} as a separate entry (4668).
\end{cousinweb}

\begin{grammarlens}[title={What you've built}]
One more word for describing what you lived and felt.
\end{grammarlens}

\subsection*{Your turn}
\begin{itemize}
\raggedright
  \item \emph{Say it:} \emph{maṇa}
  \item \emph{Say it:} \emph{maṇa}, once more
  \item \emph{Say it:} \emph{avar oru maruttuvarai maṇantār.} — He married a doctor.
  \item \emph{Recall:} say the Tamil for a dowry, then the Tamil for a blessing, then say \emph{maṇa} again
\end{itemize}

\begin{tcolorbox}[breakable,colback=teal!4,colframe=teal!35!black,title={Before you move on}]
What does \ta{மண} mean? (\textquotedblleft{}To marry\textquotedblright{}.) Say it once more.
\end{tcolorbox}
\section[\texorpdfstring{\ta{தம்பதி} (tampati)}{தம்பதி (tampati)}]{\ta{தம்பதி} (tampati) — a married couple}

\label{lesson:TA-C499-tampati}



\begin{quote}
Before the new one: say the Tamil for a blessing, then the Tamil for to marry.
\end{quote}

\subsection*{You'll want to know: \ta{தம்பதி}}
\textbf{\ta{தம்பதி}} — \emph{tampati} — \textquotedblleft{}a married couple\textquotedblright{}.

\begin{quote}
\ta{அந்தத்} \ta{தம்பதி} \ta{ஐம்பது} \ta{வருடமாக} \ta{ஒன்றாக} \ta{வாழ்கிறார்கள்}.
\end{quote}

\begin{quote}
\emph{antat tampati aimpatu varuṭamāka oṉṟāka vāḻkiṟārkaḷ.}
\end{quote}

\begin{quote}
That couple has lived together for fifty years.
\end{quote}

One word for the pair; the verb is plural.

\begin{cousinweb}
\textbf{\ta{தம்பதி}} (\emph{tampati}) is borrowed from Sanskrit \emph{dampati}, \textquotedblleft{}lord of the house\textquotedblright{}, whose dual names the married pair (Monier-Williams).
\end{cousinweb}

\begin{grammarlens}[title={What you've built}]
One more word for describing what you lived and felt.
\end{grammarlens}

\subsection*{Your turn}
\begin{itemize}
\raggedright
  \item \emph{Say it:} \emph{tampati}
  \item \emph{Say it:} \emph{tampati}, once more
  \item \emph{Say it:} \emph{antat tampati aimpatu varuṭamāka oṉṟāka vāḻkiṟārkaḷ.} — That couple has lived together for fifty years.
  \item \emph{Recall:} say the Tamil for a blessing, then the Tamil for to marry, then say \emph{tampati} again
\end{itemize}

\begin{tcolorbox}[breakable,colback=teal!4,colframe=teal!35!black,title={Before you move on}]
What does \ta{தம்பதி} mean? (\textquotedblleft{}A married couple\textquotedblright{}.) Say it once more.
\end{tcolorbox}
\section[\texorpdfstring{\ta{இல்லறம்} (illaṟam)}{இல்லறம் (illaṟam)}]{\ta{இல்லறம்} (illaṟam) — married life}

\label{lesson:TA-C499-illaram}



\begin{quote}
Before the new one: say the Tamil for to marry, then the Tamil for a married couple.
\end{quote}

\subsection*{You'll want to know: \ta{இல்லறம்}}
\textbf{\ta{இல்லறம்}} — \emph{illaṟam} — \textquotedblleft{}married life\textquotedblright{}.

\begin{quote}
\ta{அவர்களுடைய} \ta{இல்லறம்} \ta{மகிழ்ச்சியாக} \ta{இருந்தது}.
\end{quote}

\begin{quote}
\emph{avarkaḷuṭaiya illaṟam makiḻcciyāka iruntatu.}
\end{quote}

\begin{quote}
Their married life was happy.
\end{quote}

\ta{இல்} \textquotedblleft{}house\textquotedblright{} + \ta{அறம்} \textquotedblleft{}duty, virtue\textquotedblright{}.

\begin{cousinweb}
\textbf{\ta{இல்லறம்}} (\emph{illaṟam}) is a Tamil compound of native \ta{இல்} \emph{il}, \textquotedblleft{}house, home\textquotedblright{} (DEDR 494), and \ta{அறம்}, \textquotedblleft{}virtue, duty\textquotedblright{}.
\end{cousinweb}

\begin{grammarlens}[title={What you've built}]
One more word for describing what you lived and felt.
\end{grammarlens}

\subsection*{Your turn}
\begin{itemize}
\raggedright
  \item \emph{Say it:} \emph{illaṟam}
  \item \emph{Say it:} \emph{illaṟam}, once more
  \item \emph{Say it:} \emph{avarkaḷuṭaiya illaṟam makiḻcciyāka iruntatu.} — Their married life was happy.
  \item \emph{Recall:} say the Tamil for to marry, then the Tamil for a married couple, then say \emph{illaṟam} again
\end{itemize}

\begin{tcolorbox}[breakable,colback=teal!4,colframe=teal!35!black,title={Before you move on}]
What does \ta{இல்லறம்} mean? (\textquotedblleft{}Married life\textquotedblright{}.) Say it once more.
\end{tcolorbox}
